\documentclass[]{article}

\usepackage{amsmath}
\usepackage{multirow}
\usepackage{natbib}
\usepackage{graphicx} 


\begin{document}

Distinguishing between weak gravitational lensing maps generated by different hydrodynamical simulation codes presents a challenging Out-of-Distribution (OoD) detection problem, as the anomalous signal arises from subtle differences in non-Gaussian structure rather than variations in known physical parameters. We propose a two-stage framework to address this challenge. First, we employ the Wavelet Scattering Transform (WST) to extract a 217-dimensional feature vector from each map, effectively summarizing the multi-scale morphological information indicative of baryonic physics while suppressing observational noise. Second, we train a Conditional Normalizing Flow on a dataset of 25,856 simulated maps to model the probability density of these features. By conditioning the model on five cosmological and baryonic nuisance parameters, our approach learns to marginalize over known physical variations and thereby isolate true structural anomalies. During inference, a Multi-Layer Perceptron regressor estimates the conditioning parameters from a given map's WST features, and the final OoD score is computed as the negative log-likelihood under the conditional flow. On a synthetic validation set designed to test this capability, our method achieves a partial Area Under the Curve of 0.2223 in the stringent low false positive rate regime of 0.001 to 0.05, demonstrating a strong ability to identify structural deviations. The resulting scores on the test set exhibit a heavy-tailed distribution, successfully identifying a distinct population of strong OoD candidates.

\end{document}
                